After the mission command is received and before the agent starts 
carrying it out, safety rules are generated. As a result, rule 
creation is carried out independently for every job, enabling the 
framework to generate a limited set of safety attributes specific 
to the objects, activities, and risks involved in the present 
mission. The mission description serves as the main input for the 
generating process, and traces are only utilized in the verification 
stage that follows. The rules are not derived from execution traces.

To generate rules, an LLM is given controlled contextual information
about the execution environment. This context includes actions and object 
states that AI2-THOR supports, the types of embodied-agent dangers, reusable 
temporal specification patterns, and the predicate vocabulary revealed by the 
trace collection component. The temporal patterns provide common structures 
such as global avoidance, conditional response, and eventual resolution, 
while the hazard taxonomy guides the identification of potentially unsafe 
behaviours. Providing this information explicitly reduces the reliance on 
unrestricted commonsense reasoning and constrains the generated rules to 
concepts supported by the framework.

The first stage generates structured plain-English safety rules, which 
identify relevant dangers and define temporal requirements in 
machine-readable JSON format. Each rule has a rule identification, 
title, hazard category and plain-English explanation. The generated 
rules are converted into parameterized LTL templates, employing 
placeholders like \texttt{x} for relevant items. This translation 
is limited to the normalized atomic propositions defined by the trace 
vocabulary, and requirements based on unavailable information are 
ignored. Before verification, placeholders are instantiated with 
concrete object identifiers to create completely grounded LTLf formulas.